\documentclass[11pt,a4paper]{article}
\usepackage[margin=2.2cm]{geometry}
\usepackage{amsmath,amssymb,xcolor,enumitem,booktabs}
\usepackage[normalem]{ulem}
\usepackage[hidelinks]{hyperref}
\setlength{\parskip}{0.6em}\setlength{\parindent}{0pt}
\definecolor{newc}{HTML}{1f6f43}
\newcommand{\new}[1]{\textcolor{newc}{#1}}
\newcommand{\old}[1]{\textcolor{gray}{\sout{#1}}}
\newenvironment{greenblock}{\color{newc}}{}
\newenvironment{oldblock}{\color{gray}}{}
\newcommand{\lref}[1]{{\small\texttt{#1}}}
\newcommand{\moved}[1]{{\small\textcolor{gray}{\emph{[moved from #1]}}}}
\newcommand{\cmt}[1]{{\small\textcolor{gray}{\emph{[comment #1]}}}}
\newcommand{\cut}[1]{{\small\textcolor{gray}{\emph{[cut: #1]}}}}
\title{Proposed rewrite, Appendix C, Model details}
\date{25 November 2026, version 1}
\begin{document}\sloppy
\maketitle

\section*{How to read this}
This proposal settles the open points of
\texttt{appendix\_C\_modeldetails\_comments.md}. Four of them are already
done in the thesis. A1 was settled by [A109], C1 by [A141] and A4 by
[A122]. The tie order of D3 now stands in the Method
(\texttt{4.Implementation.tex}, the ties paragraph of
\lref{sec:method:storage}).

The appendix keeps only what the thesis points to. The thesis points to six
things in it. These are the sources of the model, the CRAMP comparison, the
keep-up count, the field of view at a zone edge, the formal statements and
the tie pointer. Everything else repeats the main text or is referred to
nowhere, so it goes. The appendix shrinks from four sections to three and
from about 350 lines to about 170.

Part~1 is the whole appendix as it would read. Black text stays,
\new{green} is new and \old{struck} text goes. Long cut blocks are shown as
a grey note with the reason, not printed in full. Text that only moves is
black, with a grey note saying where it came from. References are printed as
their labels. Part~2 lists the small changes outside Appendix~C. Part~3 says
what was checked and where. Part~4 holds the questions. Nothing below is in
the thesis yet.

\textbf{Labels.} Every label the chapters use stays, so no reference in the
chapters breaks. \lref{app:controllers} moves to the halo paragraph, which is
what the Method points to it for. \lref{app:measures} moves to the CRAMP
paragraph, which is what the Problem Formulation points to it for. These
labels go, and nothing in the chapters or the other appendices uses them:
\lref{app:storage}, \lref{app:storage:ties} (its one pointer is removed, see
Part~2), \lref{app:game}, \lref{app:controllers:congestion},
\lref{app:controllers:example}, \lref{tab:observationexample},
\lref{app:controllers:bound}, \lref{eq:potentialbound},
\lref{app:measures:blocked}, \lref{fig:architecturecompact} and
\lref{eq:onestepcost} (used only by the two notation tables).

% =================================================================
\section*{Part 1. The appendix as it would read}

\subsection*{C \quad Model details \hfill\lref{app:details}}

\old{The worked examples, derivations and formal statements behind the design
of the thesis, in the order of the thesis. Each keeps its number, so a
reference from the main text leads straight to it.}
\new{This appendix sets the network beside the papers it comes from, says how
the keep-up check counts, and gives the formal statements the Method gives
in words.} \hfill\cmt{D4}

\old{\textbf{C.1 Storage and inventory.}} \cut{the section opening retells
\lref{sec:pf:storage}, and C.1.1 ``Ties in the slotting costs'' only sends
the reader back to \lref{sec:method:storage}, which now holds the tie order
itself (D3)}

\old{\textbf{C.2 Controllers and the learning problem.}} \cut{the opening
retells \lref{sec:pf:controllers} and \lref{sec:pf:observations}}

\old{\textbf{C.2.1 The learning problem in full.}} \cut{the example state
repeats \lref{sec:pf:game}, the ``second reason'' has no first, and the
crowding paragraph repeats the Method (D1, D2)}

\old{\textbf{C.2.2 Which congestion values a controller reads}, first
paragraph.} \cut{says what the Method's ``Crowding'' paragraph already says,
that every controller reads only the agent's own value. The second paragraph,
on the halo, moves to C.3.3}

\old{\textbf{C.2.3 The observation of the example agent}, and Table~C.1.}
\cut{repeats the example of \lref{sec:pf:observations} ($\eta_1(v,t)=3$,
$\eta_1(w,t)=5$) and its figure (D1)}

\old{\textbf{C.2.4 The distance-label bound}, and (C.1).} \cut{nothing points
to it. Every edge costs 1, so it says only that one step shortens the
distance by at most one. Its first sentence repeats
\lref{sec:pf:observations} (B3)}

% -----------------------------------------------------------------
\subsection*{C.1 \quad The model against its sources \hfill\lref{app:controllers:sources}}

The network of \lref{sec:method:model} takes its parts from three sources.
\old{This subsection writes each in its own form and in the notation of the
thesis, so the reader can check what is the same and what differs.}
\new{Each is written below in its own form and then in the notation of the
thesis.}

\textbf{Message passing.} Gilmer et al.\ (2017) write one round of a message
passing network as
\begin{equation*}
  m_v^{(r)} = \sum_{w \in N(v)} M_r\bigl(h_v^{(r-1)}, h_w^{(r-1)}, e_{vw}\bigr),
  \qquad
  h_v^{(r)} = U_r\bigl(h_v^{(r-1)}, m_v^{(r)}\bigr),
\end{equation*}
with a message function $M_r$, an update function $U_r$ and edge features
$e_{vw}$. The message-passing step (\lref{eq:msgpass}) has the same form,
with $\phi_r$ in the role of $M_r$ and the sum, and $\psi_r$ in the role of
$U_r$. It differs in three ways. It takes the mean over the neighbours where
they take the sum. It has no edge features, which loses nothing here,
because every edge of the generated warehouses costs 1
(\lref{sec:impl:instances}). And a vertex hears its neighbours along edges in
both directions, so information also flows against one-way edges.

\textbf{Graph convolution.} Van Knippenberg et al.\ (2021) use the graph
convolution of Kipf and Welling (2017),
\begin{equation*}
  h_v^{(r)} = \sigma\Bigl( W_r \sum_{w \in N(v) \cup \{v\}}
    \frac{h_w^{(r-1)}}{\sqrt{\deg(v)\,\deg(w)}} \Bigr),
\end{equation*}
with one weight matrix $W_r$ per layer and degrees counted with the vertex
itself. \new{Their layer also reads the cost of each edge.} \old{This is a
special case of (\lref{eq:msgpass}), with fixed weights for the neighbours
and a single linear map with an activation $\sigma$ as the update.}
\new{The message-passing step has the same shape, a weight for each
neighbour and then an update. It does not hold Kipf's layer exactly, because
$\phi_r$ does not see the degrees.} The model of this thesis learns $\phi_r$
and $\psi_r$ instead, so it can weigh a neighbour by its features, for
example by whether an agent stands on it. \hfill\cmt{A2}

\old{Figure~C.1 draws the model in the form of Figure~3 of van Knippenberg et
al.\ (2021). The move score and the value head are what this model adds.
Their environment already applies the actions in an order that avoids
collisions where it can, and conflict resolution here also turns every
remaining clash into a wait.} \old{[Figure C.1, ``The model drawn as Figure~3
of (ibid.)'']} \cut{the figure shows the same blocks as the Method's
architecture figure, in symbols instead of words (E1). The value head is
named in the next paragraph's list of differences, and conflict resolution
there too. See Question~1}

\textbf{Messages between agents.} Pham et al.\ (2024) run a GNN over the
agents in an agent's field of view, with $L$ layers of graph filters,
\begin{equation*}
  h^{(l+1)}_f = \sigma_{l+1}\bigl( \textstyle\sum_{g} H_l(A)\, h^{(l)}_g \bigr),
\end{equation*}
where $f$ and $g$ number the features of a layer, not the agents. $A$ holds
the distances between the visible agents and each agent's starting state
comes from an image network over its grid window. The final states are
averaged and joined with the agent's own features. The readout
(\lref{eq:readout}) does the same with one round only. Each agent's message
$\xi_j(t)$ is its own final vertex state, the messages of the agents it hears
are averaged, and the mean is joined with its own state. One round keeps the
decentralised window at $d_{\mathrm{obs}}$ (\lref{sec:method:model}). The
communication graph (\lref{eq:commgraph}) is unweighted and directed, where
theirs weighs agents by distance.

\textbf{Crowding.} \hfill\lref{app:measures}\ \moved{C.3, second paragraph}\\
The crowding value (\lref{eq:crowding}) follows the crowd density of CRAMP
(Pham et al.\ 2024), which counts the agents in an agent's field of view and
divides by the free space there. CRAMP uses the density only in its reward,
with a threshold above which an area counts as crowded, and over a region
smaller than its field of view. Here it covers the field of view
$d_{\mathrm{obs}}$ and has no threshold. \new{It is used in three places, not
one. Every controller reads it as an input, the reward charges for it, and
its mean over the run is a reported outcome.} \old{and is averaged into a
reported outcome. That averaging is this thesis's own.} \new{The input and the
reported mean are the thesis's own.} \hfill\cmt{B2}

\textbf{Other differences from van Knippenberg et al.\ (2021).} Their agents
learn a value for each action with a deep Q-network, and here the policy is
trained with PPO. They turn an invalid move into a penalised wait, and here
the mask removes it (\lref{sec:method:model}). \new{The mask follows PRIMAL
(Sartoretti et al.\ 2019), which samples only valid moves in training
because that trained more stably than a penalty.} Their output layer gives
all move values at once, and here each move is scored by the vertex it leads
to, by the move score (\lref{eq:edgescore}). A collision ends their episode
with a penalty, and here conflict resolution turns it into a wait
(\lref{sec:method:resolution}). The step costs, $0.1$ for a move and $0.2$
for a wait, are theirs. \hfill\cmt{B1}

% -----------------------------------------------------------------
\subsection*{C.2 \quad How the keep-up check counts \hfill\lref{app:measures:keepup}}

\old{\textbf{C.3 Measures and evaluation.} The result of the thesis is the
service time gain of demand storage over fixed storage for each controller,
read as a percentage of the service time under fixed storage. A run counts
only if the fleet keeps up with the released tasks. The service time is that
of orders, from an order's release until it leaves at a dock. The service
time of put-aways and the time a controller takes to decide are reported next
to it.} \cut{retells \lref{sec:pf:measures}, and nothing below uses it (D6)}

The check (\lref{eq:throughput}) counts every task the fleet finishes in the
measured steps, also one released before them. A fleet that fell behind
earlier still works through its older tasks, and that work is what shows
whether the queue grows or holds. With 900 tasks released in the measured
steps and 600 tasks finished in them, the share is two thirds and the queue
grows by 300.

A pick counts as released when its order arrives, also while its stock is
not yet on a shelf. The order is work the fleet owes, and the backlog already
counts it. Left out, a fleet whose put-aways fail would release almost no
picks and seem to keep up. Such a fleet with 800 orders waiting for stock and
10 picks finished serves about one pick in eighty.

\old{\textbf{C.3.2 Blocked time against the shelf estimate.} The mean blocked
time $W_T$ is not the shelf estimate $\hat\omega_t(v)$. \ldots\ Neither is
derived from the other.} \cut{nothing points to it. Both are defined in the
chapters, and the difference shows in their definitions. See Question~3}

% -----------------------------------------------------------------
\subsection*{C.3 \quad Definitions in full \hfill\lref{app:defs}}

The main text states these definitions in words \new{and points here for the
formal version.} \old{and points here for the formal version. Each keeps its
number, so a reference from the main text leads straight to it.}
\hfill\cmt{D4, the same sentence as in the opening}

\textbf{C.3.1 From \lref{sec:method:storage}.} The two legs of a shelf's
trip,
\begin{equation}
  d^{\mathrm{dock}}(v) = \frac{1}{|V_{\mathrm{dock}}|} \sum_{g \in V_{\mathrm{dock}}} d_G(g, v),
  \qquad
  d^{\mathrm{work}}(k,v) = \frac{1}{|h(k)|} \sum_{w \in h(k)} d_G(v, w),
  \tag{C.1}
\end{equation}
where $h(k)$ is the category of SKU $k$, the four workstations its lines go
to. Every storage rule uses $h(k)$. The demand rules learn it from the SKU's
first order line, which always goes to a workstation of that category.
\old{A pallet that misses its home shelf under fixed storage is placed by
these legs. At 16 agents this is about 8 percent of the arriving units under
Token Passing and 40 to 60 percent under the learned controllers, whose
fleets fall behind and leave the shelves full.} \cut{run results in a
definition. The Token Passing share is already in Appendix~E
(\lref{app:load:validity}, 6 to 10 percent). See Question~2}

\textbf{C.3.2 From \lref{sec:method:model}.} \moved{after the halo, D5} The
vertex features,
\begin{equation}
  \begin{aligned}
  f(v,t) = \bigl(\, &\mathrm{onehot}(\mathrm{role}(v)),\;\; \tilde\eta_i(v,t),\;\;
                    \mathbb{1}[v \text{ occupied by another agent}],\\
                    &
                    \mathbb{1}[v \in \mathcal{V}_{\prec i}(t)],\;\; \eta_i(v,t)/50 \,\bigr) .
  \end{aligned}
  \tag{C.2}
\end{equation}

The action mask,
\begin{equation}
  \begin{aligned}
  \Pr[u_i(t) = \cdot] &= \mathrm{softmax}\bigl(\, \mathrm{logit}(\cdot) + M(v,\cdot) + M^{\prec}_i(t,\cdot) \,\bigr), \\
  M(v,\cdot) &= \begin{cases} 0 & \cdot \in \mathcal{U}(v), \\
                               -\infty & \cdot \in \mathcal{U}_i\setminus\mathcal{U}(v), \end{cases}
  \end{aligned}
  \tag{C.3}
\end{equation}
where $v = \ell_i(t)$ is the vertex agent $a_i$ stands on.

\textbf{C.3.3 From \lref{sec:method:controllers}.} \hfill\lref{app:controllers}\\
The claim mask,
\begin{equation}
  M^{\prec}_i(t,\mathrm{move\ to\ }w) =
  \begin{cases}
    -\infty & w \in \mathcal{V}_{\prec i}(t), \text{ or } w = \ell_j(t) \text{ with } j \not\prec i,\ j \neq i, \\
            & \text{or } w = \ell_j(t) \text{ with } j \prec i \text{ and } a_j \text{ has chosen } v, \\
    0       & \text{otherwise},
  \end{cases}
  \tag{C.4}
\end{equation}
with the same $v = \ell_i(t)$, so the last case is a swap.

The halo,
\begin{equation}
  H_q = \{\, v \notin Y_q : d_G(w,v) \le d_{\mathrm{obs}}
          \text{ for some } w \in Y_q \,\}.
  \tag{C.5}
\end{equation}
\moved{C.2.2, second paragraph} A section-based agent's field of view is not
cut at the edge of its zone. A reading that reaches a short way into the next
zone is still there, because the zone's window includes its halo $H_q$. The
halo exists so that the edge of a zone is not a blind spot, and cutting every
reading at the edge would undo it.

\old{\textbf{From \lref{sec:method:rl}.} The one-step cost,
$\hat{c}(v,w) = c_{\mathrm{wait}}$ if $v = w$ and $c(v,w)$ if $(v,w) \in E$.}
\cut{nothing uses it. The reward table and the reward equation of
\lref{sec:method:rl} already give $c_{\mathrm{move}} = 0.1$ and
$c_{\mathrm{wait}} = 0.2$ in words and values. Kept, it suggests a move
costs $c(v,w) = 1$ (A3)}

The order of C.3 now follows the Method. Storage comes first, then the
network (features and action mask), then choosing together (claim mask and
halo). The comment D5 asked for the halo after the features and masks.

% =================================================================
\section*{Part 2. Changes outside Appendix C}
\begin{enumerate}
  \item \texttt{4.Implementation.tex}, ties paragraph of
        \lref{sec:method:storage}. ``So each rule gives exactly one placement
        \old{(Appendix~\lref{app:storage:ties})}.'' The pointer goes with
        C.1.1, since the sentence before it already gives the tie order.
  \item \texttt{4.Implementation.tex}, the mask paragraph ``Only moves that
        exist''. ``\ldots\ So an agent can never choose a move that does not
        exist\new{, as in PRIMAL (Sartoretti et al.\ 2019) during
        training}. Van Knippenberg et al.\ instead turn such a move into a
        penalised wait, and their policy has to learn not to choose it.''
        (B1)
  \item \texttt{4.Implementation.tex}, opening of \lref{sec:method:model}.
        ``Appendix~\lref{app:controllers:sources} sets the three side by
        side\old{ and draws the model in the form of their figure}.'' Only
        if Figure~C.1 goes (Question~1).
  \item \texttt{A.Notation.tex}. The row $\hat c(v,w)$ goes (A3). The
        row $c_{\mathrm{move}}$ and $c_{\mathrm{wait}}$ stay.
  \item \texttt{styles/eqnames.tex}. The names of \lref{eq:potentialbound} and
        \lref{eq:onestepcost} go. Then \texttt{tools/sync\_xrefs.py --build}.
  \item \texttt{GAPS.md}. One entry, [A143], text only, no model change.
        Every changed line tagged \texttt{\% [A143]}.
\end{enumerate}

% =================================================================
\section*{Part 3. What was checked}
\begin{itemize}[nosep]
  \item \textbf{Which labels the thesis uses.} A search of every
        \texttt{\textbackslash ref}, \texttt{\textbackslash eqref} and
        \texttt{\textbackslash nref} in \texttt{Chapters/} and
        \texttt{Appendices/}. The chapters point to
        \lref{app:storage:ties}, \lref{app:controllers},
        \lref{app:controllers:sources}, \lref{app:measures},
        \lref{app:measures:keepup}, and to the equations
        \lref{eq:tripdistance}, \lref{eq:halo}, \lref{eq:features},
        \lref{eq:mask} and \lref{eq:claimmask}. Nothing else in
        Appendix~C is referred to.
  \item \textbf{PRIMAL and the mask.} Sartoretti et al.\ (2019), Section
        III-B: ``During training, actions are sampled only from valid
        actions \ldots\ this approach enables more stable training, compared
        to giving negative rewards''. At test time an invalid move becomes a
        wait (\texttt{papers/1809.03531v3.pdf}).
  \item \textbf{Edge attributes of van Knippenberg et al.} Their Section 2.3:
        a GCN embeds the graph ``including any node and edge attributes'',
        and their graphs have both (\texttt{papers/knippenberg21a.pdf}).
  \item \textbf{Crowding in three places.} The Method's ``Crowding''
        paragraph (input of every controller) and the reward table, row
        crowding, weight $n_{\mathrm{crowd}} = 0.1$
        (\texttt{4.Implementation.tex}).
  \item \textbf{The step costs.} The reward table gives $c_{\mathrm{move}}
        = 0.1$ and $c_{\mathrm{wait}} = 0.2$ from van Knippenberg et al.,
        Table~1. The code charges the same (\texttt{reward\_function.py},
        \texttt{c\_move}). No code uses a one-step cost $\hat c$.
  \item \textbf{The two figures.} \texttt{Images/architecture.pdf} (Method)
        and \texttt{Images/architecture\_compact.pdf} (Figure~C.1) hold the
        same blocks. The Method's figure has one agent's network, the agents,
        conflict resolution and the warehouse, in words. Figure~C.1 has the
        same in symbols ($L$, $u_i$, $R_i$).
  \item \textbf{The overflow numbers.} ``About 8 percent under Token Passing,
        40 to 60 percent under the learned controllers'' come from [A122]
        (23 November). I found no rows for the learned share in
        \texttt{results/}. Appendix~E gives 6 to 10 percent for Token
        Passing from the updated load run.
\end{itemize}

% =================================================================
\section*{Part 4. Questions}
\begin{enumerate}
  \item \textbf{Figure C.1.} It shows the same network as the Method's
        figure, in symbols. The draft drops it. The other choice is to keep
        it with the caption ``The model drawn as in Figure~3 of van
        Knippenberg et al.\ (2021).''
        \\ \hspace*{1em}$\square$ drop \quad $\square$ keep, new caption \quad Note:
  \item \textbf{The learned overflow share.} The draft takes it out of
        Appendix~C. Should the share of arriving stock that misses its home
        shelf under the learned controllers go into Appendix~L (``Stock in
        the test runs''), recounted from the long 16-agent runs? Or is it
        left out?
        \\ \hspace*{1em}$\square$ recount, add to L \quad $\square$ leave out \quad Note:
  \item \textbf{Blocked time against the shelf estimate.} Nothing points
        to it, so the draft cuts it. The other choice is one sentence where
        $W_T$ is defined in \lref{sec:pf:measures}: ``$W_T$ looks back over a
        run, and the shelf estimate is a forecast that steers storage.'' That
        would touch the Problem Formulation.
        \\ \hspace*{1em}$\square$ cut \quad $\square$ one sentence in the PF \quad Note:
  \item \textbf{Where the keep-up count sits.} The draft keeps it in C.2.
        It could move to Appendix~E, next to ``The keep-up threshold''
        (\lref{app:cal:load}), where the threshold is set. The label stays,
        so the Problem Formulation's pointer still works.
        \\ \hspace*{1em}$\square$ keep in C \quad $\square$ move to E \quad Note:
  \item \textbf{Files nothing includes.} \texttt{Chapters/2.Notation.tex} is
        an older copy of \texttt{A.Notation.tex}, and no file includes it.
        The same goes for \texttt{tmp1.tex}, \texttt{5.Experiments.before.*},
        \texttt{Project.tex}, \texttt{Thesis\_idea.tex},
        \texttt{Timeplan.tex} and \texttt{Helper.tex}. Delete them from the
        Overleaf project (they stay in its history)? This is separate from
        Appendix~C.
        \\ \hspace*{1em}$\square$ delete \quad $\square$ keep \quad Note:
\end{enumerate}

\end{document}
